% generated by rb.py note; do not edit
\newcommand{\ConcDecls}{4}
\newcommand{\Pin}{72fead10}
\newcommand{\PinText}{that of commit {\leanfont 72fead10}}
\newcommand{\ArtifactName}{grh-compact}
\expandafter\newcommand\csname LN@rh\endcsname{\href{https://github.com/lsernau/grh-lean/blob/72fead102f663d5c7a407964a13dc0cf9d977207/grh-compact/src/Statements.lean\#L3-L4}{{\leanfont RiemannZetaZerosOnLine}}}
\expandafter\newcommand\csname LQ@rh\endcsname{\leanquoteline{0}{def RiemannZetaZerosOnLine : Prop :=}\leanquoteline{2}{∀ s : ℂ, riemannZeta s = 0 → (¬∃ n : ℕ, s = -2 * ((n : ℂ) + 1)) → s ≠ 1 → s.re = 1 / 2}}
\expandafter\newcommand\csname LN@halfPlane\endcsname{\href{https://github.com/lsernau/grh-lean/blob/72fead102f663d5c7a407964a13dc0cf9d977207/grh-compact/src/Statements.lean\#L6-L8}{{\leanfont DirichletLNonvanishingRightOfHalf}}}
\expandafter\newcommand\csname LQ@halfPlane\endcsname{\leanquoteline{0}{def DirichletLNonvanishingRightOfHalf : Prop :=}\leanquoteline{2}{∀ (N : ℕ) [NeZero N] (χ : DirichletCharacter ℂ N) (s : ℂ),}\leanquoteline{4}{1 / 2 < s.re → (χ ≠ 1 ∨ s ≠ 1) → DirichletCharacter.LFunction χ s ≠ 0}}
\expandafter\newcommand\csname LN@strip\endcsname{\href{https://github.com/lsernau/grh-lean/blob/72fead102f663d5c7a407964a13dc0cf9d977207/grh-compact/src/Statements.lean\#L10-L12}{{\leanfont DirichletLZerosOnLine}}}
\expandafter\newcommand\csname LQ@strip\endcsname{\leanquoteline{0}{def DirichletLZerosOnLine : Prop :=}\leanquoteline{2}{∀ (N : ℕ) [NeZero N] (χ : DirichletCharacter ℂ N) (s : ℂ),}\leanquoteline{4}{0 < s.re → s.re < 1 → DirichletCharacter.LFunction χ s = 0 → s.re = 1 / 2}}
\expandafter\newcommand\csname LN@stripPrim\endcsname{\href{https://github.com/lsernau/grh-lean/blob/72fead102f663d5c7a407964a13dc0cf9d977207/grh-compact/src/Statements.lean\#L14-L16}{{\leanfont PrimitiveDirichletLZerosOnLine}}}
\expandafter\newcommand\csname LQ@stripPrim\endcsname{\leanquoteline{0}{def PrimitiveDirichletLZerosOnLine : Prop :=}\leanquoteline{2}{∀ (N : ℕ) [NeZero N] (χ : DirichletCharacter ℂ N), χ.IsPrimitive →}\leanquoteline{4}{∀ s : ℂ, 0 < s.re → s.re < 1 → DirichletCharacter.LFunction χ s = 0 → s.re = 1 / 2}}
\newcommand{\ConcordanceRows}{\crow{rh}{\href{https://github.com/lsernau/grh-lean/blob/72fead102f663d5c7a407964a13dc0cf9d977207/grh-compact/src/Statements.lean\#L3-L4}{{\leanfont RiemannZetaZerosOnLine}}}{Statements.lean}{3--4}
\crow{halfPlane}{\href{https://github.com/lsernau/grh-lean/blob/72fead102f663d5c7a407964a13dc0cf9d977207/grh-compact/src/Statements.lean\#L6-L8}{{\leanfont DirichletLNonvanishingRightOfHalf}}}{Statements.lean}{6--8}
\crow{strip}{\href{https://github.com/lsernau/grh-lean/blob/72fead102f663d5c7a407964a13dc0cf9d977207/grh-compact/src/Statements.lean\#L10-L12}{{\leanfont DirichletLZerosOnLine}}}{Statements.lean}{10--12}
\crow{stripPrim}{\href{https://github.com/lsernau/grh-lean/blob/72fead102f663d5c7a407964a13dc0cf9d977207/grh-compact/src/Statements.lean\#L14-L16}{{\leanfont PrimitiveDirichletLZerosOnLine}}}{Statements.lean}{14--16}}
% the notations of the defined functions: one slot for each explicit argument of the Lean declaration
\newcommand{\UnitMembers}{4}
\newcommand{\UnitInNote}{4}
\newcommand{\UnitProofs}{0}
\newcommand{\UnitSteps}{0}
\newcommand{\UnitAllSteps}{0}
